\section{Physical propagation on the explicit grid}
\label{sec:huang_physical_propagation}

The following construction applies to the original nearest-neighbor interactions.
It supplies the full-chain propagator approximation used in
\cite[Section 5.4, Lemma 19]{Huang2015GroundState}, following
\cite[Figure 2 and equations (2)--(13)]{Osborne2006Dynamics}.

\begin{definition}[Within-block and omitted interactions]
    \label{def:huang_grid_hamiltonian}
    \lean{TNLean.MPS.AreaLaw.chainGridBoundaryEdge,
        TNLean.MPS.AreaLaw.chainGridBoundaryEdges,
        TNLean.MPS.AreaLaw.chainGridEdgeOwner,
        TNLean.MPS.AreaLaw.chainGridBlockHamiltonian,
        TNLean.MPS.AreaLaw.chainGridRetainedHamiltonian,
        TNLean.MPS.AreaLaw.chainGridBoundaryTerm}
    \leanok
    \uses{def:huang_chain_grid}
    Write $H=\sum_{i=0}^{L-2}h_i$, where $h_i$ acts on sites $i,i+1$.
    At a grid boundary $c_k$, let $e_k=c_k-1$ be the omitted edge index,
    and put $E=\{e_k\}_k$. The owner of an edge $i$ is $j(i)=\lfloor i/w\rfloor$.
    Define
    \[
        A_j=\sum_{i\notin E,\ j(i)=j}h_i,\qquad
        A=\sum_jA_j,\qquad B_k=-h_{e_k}.
    \]
\end{definition}

\begin{lemma}[Exact Hamiltonian decomposition]
    \label{thm:huang_grid_hamiltonian_decomposition}
    \lean{TNLean.MPS.AreaLaw.chainGridBoundaryEdge_val_succ,
        TNLean.MPS.AreaLaw.chainGridBoundaryEdge_injective,
        TNLean.MPS.AreaLaw.chainGridEdge_subset_ownerBlock,
        TNLean.MPS.AreaLaw.chainGridRetainedHamiltonian_eq_sum_filter,
        TNLean.MPS.AreaLaw.sum_eq_chainGridRetained_sub_boundary,
        TNLean.MPS.AreaLaw.chainGridBlockHamiltonian_mem_supported,
        TNLean.MPS.AreaLaw.isHermitian_chainGridBlockHamiltonian}
    \leanok
    \uses{def:huang_grid_hamiltonian, thm:huang_chain_grid_edges}
    The omitted edge indices are distinct and satisfy $e_k+1=c_k$.
    Every non-omitted edge is contained in its assigned block, and
    \[
        A=\sum_{i\notin E}h_i,\qquad H=A-\sum_kB_k.
    \]
    Each $A_j$ is supported on $C_j$. If the original terms are Hermitian,
    then every $A_j$ is Hermitian.
\end{lemma}
\begin{proof}
    \leanok
    The explicit edge classification separates the original edge indices into
    internal and boundary edges. Group the internal terms by their unique owners.
    Distinct boundaries give distinct omitted edges. Support and Hermiticity are
    preserved under finite sums.
\end{proof}

\begin{lemma}[Exact evolution of the retained interaction]
    \label{thm:huang_grid_exact_evolution}
    \lean{TNLean.MPS.AreaLaw.hermitianUnitaryPath_mem_supported,
        TNLean.MPS.AreaLaw.commute_chainGridBlockHamiltonian,
        TNLean.MPS.AreaLaw.hermitianUnitaryPath_chainGridRetained_eq_prod,
        TNLean.MPS.AreaLaw.chainGridBlockEvolution_mem_supported}
    \leanok
    \uses{thm:huang_grid_hamiltonian_decomposition}
    The block Hamiltonians commute, and, in their prescribed order,
    \[
        e^{itA}=\prod_j e^{itA_j}.
    \]
    Each factor is supported on its actual block. More generally, the exponential
    path of a supported operator remains supported on the same region.
\end{lemma}
\begin{proof}
    \leanok
    Operators with disjoint supports commute. The exponential of a finite sum of
    commuting operators factors into their ordered exponentials. An exponential
    commutes with every operator commuting with its generator; the complementary
    commutant characterization gives the support assertion.
\end{proof}

\begin{lemma}[Derived localization of the omitted interactions]
    \label{thm:huang_grid_generator_localization}
    \lean{TNLean.MPS.AreaLaw.nearestNeighborChainSupport,
        TNLean.MPS.AreaLaw.chainGridRetainedTerm,
        TNLean.MPS.AreaLaw.sum_chainGridRetainedTerm,
        TNLean.MPS.AreaLaw.card_nearestNeighborChainSupport_containing_le_two,
        TNLean.MPS.AreaLaw.chainGridRetained_site_budget,
        TNLean.MPS.AreaLaw.chainGridBoundaryTerm_supported,
        TNLean.MPS.AreaLaw.chainGridLRVelocity,
        TNLean.MPS.AreaLaw.chainGridLRPrefactor,
        TNLean.MPS.AreaLaw.chainGridLRPrefactor_pos,
        TNLean.MPS.AreaLaw.norm_chainGridBoundaryTerm_sub_siteExpectation_le}
    \leanok
    \uses{thm:huang_grid_hamiltonian_decomposition, thm:huang_chain_grid_buffers}
    Suppose the original terms are Hermitian and $\|h_i\|\le1$. Set
    \[
        v=8e^2+1,\qquad C=\frac{8(2+8e^2)e}{1-e^{-1/2}}>0.
    \]
    The retained original terms are obtained by replacing precisely the omitted
    terms by zero. At most two designated supports meet any site, so their site
    norm budget is at most two. If $\mathbb E_{K_k}$ is the normalized site expectation
    onto the actual radius-$r$ buffer, then
    \[
        \|e^{isA}B_ke^{-isA}-\mathbb E_{K_k}(e^{isA}B_ke^{-isA})\|
        \le C e^{v|s|-r/2}.
    \]
\end{lemma}
\begin{proof}
    \leanok
    Apply the chain Lieb--Robinson localization bound to the retained original
    nearest-neighbor terms. The site budget is two, the omitted term has norm at
    most one, and its designated pair has at most two sites. The exact neighborhood
    identity identifies the localization region with $K_k$.
\end{proof}

\begin{lemma}[One physical two-layer approximation at every cut]
    \label{thm:huang_grid_propagator}
    \lean{TNLean.MPS.AreaLaw.exists_chainGrid_propagator_approximation}
    \leanok
    \uses{thm:huang_grid_generator_localization, thm:huang_grid_exact_evolution,
        thm:huang_two_interval_layer_ranks}
    Let $2r+2\le w$ and let $m=\lfloor(L-1)/w\rfloor$. The original terms determine
    a family $K(t)$ such that
    \[
        \|e^{-itH}-K(t)\|\le mC|t|e^{v|t|-r/2},\qquad
        \operatorname{OSR}_k(K(t))\le q^{2(w+2r+2)}
    \]
    for every real $t$ and every cut $k$. The same operator $K(t)$ occurs at all cuts.
\end{lemma}
\begin{proof}
    \leanok
    Construct the projected relative-evolution paths on the disjoint boundary
    buffers. Their ordered product approximates the full relative evolution with
    error bounded by $|t|$ times the sum of the generator errors. Multiply by the
    exact retained evolution, written as its actual block product. This is a product
    of two disjoint interval layers of widths $w$ and $2r+2$, giving the rank bound.
\end{proof}

\begin{lemma}[Explicit radius at prescribed accuracy]
    \label{thm:huang_propagation_radius}
    \lean{TNLean.MPS.AreaLaw.chainPropagationRadius,
        TNLean.MPS.AreaLaw.chainPropagationRadius_error_le,
        TNLean.MPS.AreaLaw.chainPropagationRadius_width_le}
    \leanok
    Let $T\ge0$, $\varepsilon>0$, $\delta=\min\{\varepsilon,1\}$, and
    \[
        r=\left\lceil2\left(vT+
        \log\frac{(L+1)C(T+1)}{\delta}\right)\right\rceil_+.
    \]
    For every integer $m\le L$, one has $mCTe^{vT-r/2}\le\varepsilon$. Moreover,
    \[
        2r+2\le4\bigl((v+1)T+\log(L+1)+|\log C|+\log(1/\delta)\bigr)+4.
    \]
\end{lemma}
\begin{proof}
    \leanok
    The ceiling bound gives $e^{vT-r/2}\le\delta/((L+1)C(T+1))$.
    For the width bound, expand the logarithm of the product and use
    $\log(T+1)\le T$ and the elementary upper bound on the natural ceiling.
\end{proof}

\begin{lemma}[Uniform physical propagation at prescribed precision]
    \label{thm:huang_physical_propagation}
    \lean{TNLean.MPS.AreaLaw.exists_physical_chain_propagator_approximation,
        TNLean.MPS.AreaLaw.exists_positive_chain_propagator_approximation}
    \leanok
    \uses{thm:huang_grid_propagator, thm:huang_propagation_radius}
    Let $q\ge1$, let $H=\sum_ih_i$ be the original nearest-neighbor Hamiltonian,
    and suppose $0\le h_i\le I$. For $T\ge0$ and $\varepsilon>0$, choose the radius
    above and put $w=2r+2$. There is one family $K(t)$ with
    \[
        \|e^{-itH}-K(t)\|\le\varepsilon\quad (|t|\le T),\qquad
        \operatorname{OSR}_k(K(t))\le q^{8r+8}\quad(t\in\mathbb R, k\ge0).
    \]
    The same conclusion holds for arbitrary Hermitian nearest-neighbor terms
    with norms at most one. No partition, supported propagator or localization
    error is assumed as additional input.
\end{lemma}
\begin{proof}
    \leanok
    Positive contractions are Hermitian and have norm at most one. Apply the
    constructed grid approximation. The error bound increases with $|t|$ on
    $[0,T]$, and there are at most $L$ boundary labels. The radius estimate gives
    the desired error; substituting $w=2r+2$ gives the rank exponent $8r+8$.
\end{proof}
